\documentclass[10pt,a4paper]{amsart}
\usepackage{amsmath,amssymb,amsthm,graphicx}
% For commutative diagrams we use Xy-pic
\usepackage[all]{xy}

% We use 2cell for 2-commutative diagrams.
\xyoption{2cell}
% Theorem environments.
%
\theoremstyle{plain}
\newtheorem{theorem}[subsection]{Theorem}
\newtheorem{proposition}[subsection]{Proposition}
\newtheorem{lemma}[subsection]{Lemma}

\theoremstyle{definition}
\newtheorem{definition}[subsection]{Definition}
\newtheorem{example}[subsection]{Example}
\newtheorem{exercise}[subsection]{Exercise}
\newtheorem{situation}[subsection]{Situation}

\theoremstyle{remark}
\newtheorem{remark}[subsection]{Remark}
\newtheorem{remarks}[subsection]{Remarks}

\numberwithin{equation}{subsection}

% Macros
%
\def\lim{\mathop{\mathrm{lim}}\nolimits}
\def\colim{\mathop{\mathrm{colim}}\nolimits}
\def\Spec{\mathop{\mathrm{Spec}}}
\def\Hom{\mathop{\mathrm{Hom}}\nolimits}
\def\Ext{\mathop{\mathrm{Ext}}\nolimits}
\def\SheafHom{\mathop{\mathcal{H}\!\mathit{om}}\nolimits}
\def\SheafExt{\mathop{\mathcal{E}\!\mathit{xt}}\nolimits}
\def\Sch{\mathit{Sch}}
\def\Mor{\mathop{\mathrm{Mor}}\nolimits}
\def\Ob{\mathop{\mathrm{Ob}}\nolimits}
\def\Sh{\mathop{\mathit{Sh}}\nolimits}
\def\NL{\mathop{N\!L}\nolimits}
\def\CH{\mathop{\mathrm{CH}}\nolimits}
\def\proetale{{pro\text{-}\acute{e}tale}}
\def\etale{{\acute{e}tale}}
\def\QCoh{\mathit{QCoh}}
\def\Ker{\mathop{\mathrm{Ker}}}
\def\Im{\mathop{\mathrm{Im}}}
\def\Coker{\mathop{\mathrm{Coker}}}
\def\Coim{\mathop{\mathrm{Coim}}}

% Boxtimes
%
\DeclareMathSymbol{\boxtimes}{\mathbin}{AMSa}{"02}

%
% Macros for moduli stacks/spaces
%
\def\QCohstack{\mathcal{QC}\!\mathit{oh}}
\def\Cohstack{\mathcal{C}\!\mathit{oh}}
\def\Spacesstack{\mathcal{S}\!\mathit{paces}}
\def\Quotfunctor{\mathrm{Quot}}
\def\Hilbfunctor{\mathrm{Hilb}}
\def\Curvesstack{\mathcal{C}\!\mathit{urves}}
\def\Polarizedstack{\mathcal{P}\!\mathit{olarized}}
\def\Complexesstack{\mathcal{C}\!\mathit{omplexes}}
% \Pic is the operator that assigns to X its picard group, usage \Pic(X)
% \Picardstack_{X/B} denotes the Picard stack of X over B
% \Picardfunctor_{X/B} denotes the Picard functor of X over B
\def\Pic{\mathop{\mathrm{Pic}}\nolimits}
\def\Picardstack{\mathcal{P}\!\mathit{ic}}
\def\Picardfunctor{\mathrm{Pic}}
\def\Deformationcategory{\mathcal{D}\!\mathit{ef}}
\newenvironment{reference}{}{}
\newenvironment{slogan}{}{}
\newenvironment{history}{}{}
\usepackage{cleveref}
\crefname{originalsection}{original section}{original sections}
\Crefname{originalsection}{Original section}{Original sections}
\makeatletter
\newlabel{situation-regular-model}{{0C61}{source}{}{sourcepointer-1}{}}
\newlabel{situation-regular-model@cref}{{[originalsection][0][]0C61}{[1][0][]source}{}{}{}}
\makeatother
\crefname{originallemma}{original lemma}{original lemmas}
\Crefname{originallemma}{Original lemma}{Original lemmas}
\makeatletter
\newlabel{more-morphisms-lemma-torsion-pic-thickening}{{0C6S}{source}{}{sourcepointer-2}{}}
\newlabel{more-morphisms-lemma-torsion-pic-thickening@cref}{{[originallemma][0][]0C6S}{[1][0][]source}{}{}{}}
\makeatother
\crefname{originallemma}{original lemma}{original lemmas}
\Crefname{originallemma}{Original lemma}{Original lemmas}
\makeatletter
\newlabel{coherent-lemma-power-ideal-kills-sheaf}{{01Y9}{source}{}{sourcepointer-3}{}}
\newlabel{coherent-lemma-power-ideal-kills-sheaf@cref}{{[originallemma][0][]01Y9}{[1][0][]source}{}{}{}}
\makeatother
\crefname{originaltheorem}{original theorem}{original theorems}
\Crefname{originaltheorem}{Original theorem}{Original theorems}
\makeatletter
\newlabel{coherent-theorem-formal-functions}{{02OC}{source}{}{sourcepointer-4}{}}
\newlabel{coherent-theorem-formal-functions@cref}{{[originaltheorem][0][]02OC}{[1][0][]source}{}{}{}}
\makeatother
\crefname{originallemma}{original lemma}{original lemmas}
\Crefname{originallemma}{Original lemma}{Original lemmas}
\makeatletter
\newlabel{more-algebra-lemma-isomorphic-completions}{{09BC}{source}{}{sourcepointer-5}{}}
\newlabel{more-algebra-lemma-isomorphic-completions@cref}{{[originallemma][0][]09BC}{[1][0][]source}{}{}{}}
\makeatother
\title{Torsion of order prime to the residue characteristic embeds from the Picard group of the model into that of the reduction of the special fibre}
\author{The Stacks Project Authors}
\date{stacks-project:models.tex:local-snapshot-20261012; GNU FDL 1.2 or later}
\begin{document}
\maketitle
\noindent\textit{Editorial note.} The counted unit is the Stacks Project Lemma lemma-torsion-embeds at source lines 5994-6002 together with its proof at source lines 6004-6042, a straight adjacent pairing in the chapter file. The chapter is the upstream file models.tex; the pool directory carries the longer name stacks-models and is not used as the chapter name anywhere in this card.\par\medskip
\noindent\textit{Editorial note.} This card has no context span; the situation named by the statement is declared as a pointer with the kind section, because the assembler's whitelist of kinds has no entry for a situation.\par\medskip
\noindent\textit{Editorial note.} The retained proof quotes four results of the More on Morphisms, Cohomology of Schemes and More on Algebra chapters; all four are declared as non-load-bearing navigation pointers with their stable Stacks tags looked up in the shared tag index, and none of their statements is reproduced inside the span.\par\medskip
\noindent\textit{Editorial note.} The retained proof span starts at line 6004, the line of the opening marker read from the source, and contains exactly one proof environment.\par\medskip
\noindent\textit{Editorial note.} No citation occurs in the retained spans, so no bibliography entry is supplied.\par\medskip
\noindent\textit{Allowed standard prerequisites (editorial).}\begin{itemize}
\item The regular-model situation, torsion in Picard groups and the thickenings of the special fibre. The situation fixes a discrete valuation ring with residue field, a smooth projective curve over its fraction field, a regular proper model and the components of the special fibre; the reduction of the fibre sits inside it as a closed subscheme whose successive thickenings are cut out by powers of the uniformizer, and the whole model is recovered as the formal completion along the reduction. For a class in the Picard group that is annihilated by an integer prime to the residue characteristic, triviality on the reduction implies triviality on every thickening and then on the completion, and the completion determines the scheme in the relevant sense. The author's own lemmas of the chapter and the quoted results of the Coherent, More on Algebra and More on Morphisms chapters are the allowed inputs.. The counted lemma is the comparison that lets the chapter move torsion information between the model and the reduction of its special fibre, and it is what makes the study of the Picard group of a model effective.
\end{itemize}\medskip
\ifdefined\hypertarget\hypertarget{sourcepointer-1}{}\fi
\noindent\textit{Source reference.} Original section 0C61: focused-seeds/AGP/sources/stacks-snapshot-20261011/tags.txt, source lines 15281--15281. Pointer to the situation named by the retained statement, declared with the kind section because the assembler's whitelist of kinds has no entry for a situation. The number is the stable Stacks tag of the target, looked up in the shared tag index.\par\medskip
\ifdefined\hypertarget\hypertarget{sourcepointer-2}{}\fi
\noindent\textit{Source reference.} Original lemma 0C6S: focused-seeds/AGP/sources/stacks-snapshot-20261011/tags.txt, source lines 15307--15307. Allowed input of the More on Morphisms chapter quoted by the retained proof for the comparison of Picard groups under a thickening, which is what propagates triviality from the reduction to the thickenings; the pointer is registered so the reference resolves and the cited statement is not reproduced inside the span. The number is the stable Stacks tag of the target, looked up in the shared tag index.\par\medskip
\ifdefined\hypertarget\hypertarget{sourcepointer-3}{}\fi
\noindent\textit{Source reference.} Original lemma 01Y9: focused-seeds/AGP/sources/stacks-snapshot-20261011/tags.txt, source lines 2548--2548. Allowed input of the Cohomology of Schemes chapter quoted by the retained proof to see that a sheaf supported on the special fibre is killed by a power of the uniformizer; the pointer is registered so the reference resolves and the cited statement is not reproduced inside the span. The number is the stable Stacks tag of the target, looked up in the shared tag index.\par\medskip
\ifdefined\hypertarget\hypertarget{sourcepointer-4}{}\fi
\noindent\textit{Source reference.} Original theorem 02OC: focused-seeds/AGP/sources/stacks-snapshot-20261011/tags.txt, source lines 3511--3511. Allowed input of the Cohomology of Schemes chapter quoted by the retained proof as the formal-functions theorem with which the passage to the completion is made; the pointer is registered so the reference resolves and the cited theorem is not reproduced inside the span. The number is the stable Stacks tag of the target, looked up in the shared tag index.\par\medskip
\ifdefined\hypertarget\hypertarget{sourcepointer-5}{}\fi
\noindent\textit{Source reference.} Original lemma 09BC: focused-seeds/AGP/sources/stacks-snapshot-20261011/tags.txt, source lines 11790--11790. Allowed input of the More on Algebra chapter quoted by the retained proof for the comparison of completions, which is how the conclusion is brought back to the model; the pointer is registered so the reference resolves and the cited statement is not reproduced inside the span. The number is the stable Stacks tag of the target, looked up in the shared tag index.\par\medskip
\begin{lemma}
\label{lemma-torsion-embeds}
In Situation \ref{situation-regular-model} let $h$ be an integer
prime to the characteristic of $k$. Then the map
$$
\Pic(X)[h] \longrightarrow \Pic((X_k)_{red})[h]
$$
is injective.
\end{lemma}

\begin{proof}
Observe that $X \times_{\Spec(R)} \Spec(R/\pi^n)$ is a finite
order thickening of $(X_k)_{red}$ (this follows for example from
Cohomology of Schemes, Lemma \ref{coherent-lemma-power-ideal-kills-sheaf}).
Thus the canonical map
$\Pic(X \times_{\Spec(R)} \Spec(R/\pi^n)) \to \Pic((X_k)_{red})$
identifies $h$ torsion by
More on Morphisms, Lemma \ref{more-morphisms-lemma-torsion-pic-thickening}
and our assumption on $h$.
Thus if $\mathcal{L}$ is an $h$-torsion invertible sheaf on $X$
which restricts to the trivial sheaf on $(X_k)_{red}$ then
$\mathcal{L}$ restricts to the trivial sheaf on
$X \times_{\Spec(R)} \Spec(R/\pi^n)$ for all $n$.
We find
\begin{align*}
H^0(X, \mathcal{L})^\wedge
& =
\lim H^0(X \times_{\Spec(R)} \Spec(R/\pi^n),
\mathcal{L}|_{X \times_{\Spec(R)} \Spec(R/\pi^n)}) \\
& \cong
\lim H^0(X \times_{\Spec(R)} \Spec(R/\pi^n),
\mathcal{O}_{X \times_{\Spec(R)} \Spec(R/\pi^n)}) \\
& =
R^\wedge
\end{align*}
using the theorem on formal functions
(Cohomology of Schemes, Theorem \ref{coherent-theorem-formal-functions})
for the first and last equality and for example
More on Algebra, Lemma \ref{more-algebra-lemma-isomorphic-completions}
for the middle isomorphism. Since $H^0(X, \mathcal{L})$ is a finite
$R$-module and $R$ is a discrete valuation ring, this means that
$H^0(X, \mathcal{L})$ is free of rank $1$ as an $R$-module.
Let $s \in H^0(X, \mathcal{L})$ be a basis element.
Then tracing back through the isomorphisms above we see
that $s|_{X \times_{\Spec(R)} \Spec(R/\pi^n)}$ is a trivialization
for all $n$. Since the vanishing locus of $s$ is closed in $X$
and $X \to \Spec(R)$ is proper we conclude that the vanishing
locus of $s$ is empty as desired.
\end{proof}

\end{document}
